\begingroup
\renewcommand{\arraystretch}{1.45}
\fancyhead[L]{\small ZUR THEORIE DER NEUEN TRANSZENDENTEN. V.}
\providecommand{\GaussCircleArrow}{\ensuremath{\mathord{\vcenter{\hbox{\setlength{\unitlength}{1pt}\begin{picture}(10,11)\linethickness{.55pt}\put(3.3,3.3){\circle{5.2}}\put(5.1,5.1){\line(1,1){3.2}}\put(5.6,8.3){\line(1,0){2.7}}\put(8.3,5.6){\line(0,1){2.7}}\end{picture}}}}}}

\Needspace{22\baselineskip}
\par\medskip\Needspace{5\baselineskip}\pdfbookmark[1]{Art. 9}{transscendentes-v-art-9}\centerline{[9.]}\nobreak\smallskip
Bei der $5^{\mathrm{plic.}}$ gilt (Aug. 29)
\[
\begin{aligned}
pp-PP&= A\,.\,\surd\frac{p}{P}, &5PP-pp&=B\surd\frac{P}{p}\\[8pt]
qq-QQ&=-A\,.\,\surd\frac{q}{Q}, &5QQ-qq&=B\surd\frac{Q}{q}\\[8pt]
\GaussCircleArrow\quad rr-RR&=A\,.\,\surd\frac{r}{R},
&5RR-rr&=-B\surd\frac{R}{r}
\end{aligned}
\]

\Needspace{25\baselineskip}
\[
\begin{gathered}AB=-2ppPP+2qqQQ+2rrRR=[16pqrPQR]^{\frac23},
\\[8pt]\quad{}\qquad\frac{A}{B}=\frac{PQR}{pqr}\end{gathered}
\]
\[
A=2^{\frac43}\,.\,\frac{(PQR)^{\frac56}}{(pqr)^{\frac16}},
\qquad B=2^{\frac43}\,.\,\frac{(pqr)^{\frac56}}{(PQR)^{\frac16}}
\]
also
\[
pp=(QR+qr)\,.\,\mathord{{}^{6}\!\surd}\frac{ppPP}{16qrQR}\quad\text{usw.}
\]
\[
5=-\frac{qr}{QR}+\frac{pr}{PR}+\frac{pq}{PQ}
\]
Siehe weiter unten [Art. 11.]


\Needspace{22\baselineskip}
Für die Ableitung dient u. a.
\[
\begin{gathered}
P(x^5,x^4)\,.\,P(x^5,-xx)
=P(x^{10},-xx)\,.\,P(x^{10},-x^6)\\[8pt]\quad{}+xP(x^{10},-x^4)\,.\,P(x^{10},-x^8)\\[8pt]
P(x^5,-x^4)\,.\,P(x^5,xx)
=P(x^{10},-xx)\,.\,P(x^{10},-x^6)\\[8pt]\quad{}-xP(x^{10},-x^4)\,.\,
P(x^{10},-x^8)
\end{gathered}
\]


\Needspace{24\baselineskip}
\SourceLeaf{485}
also durch Multiplikation
\[
\begin{gathered}
Q(x^{10},1)^2\, .\,P(x^{10},-x^8)\, .\,P(x^{10},-x^4)\\[8pt]
\quad=P(x^{10},-xx)^2\, .\,P(x^{10},-x^6)^2\\[8pt]\quad{}-xxP(x^{10},-x^4)^2\, .\,P(x^{10},-x^8)^2
\end{gathered}
\]
\[
\frac{(Fx^{10})^4}{(Fx^{20})^2}
=\frac{(Fx^4)^3(Fx^{10})^2}{Fx^2\, .\,Fx^{20}}
-xx\frac{Fx^2\, .\,(Fx^{20})^3}{Fx^4\, .\,Fx^{10}}
\]

\Needspace{13\baselineskip}
was mit \GaussCircleArrow{} identisch ist, nämlich
\[
p^{\frac16}q^{\frac16}(\tfrac12r)^{\frac26}
P^{\frac{1}{3}}
Q^{\frac{1}{3}}
(\tfrac12R)^{\frac23}=
\]

\Needspace{28\baselineskip}
Auch beweist man leicht
\[
\begin{gathered}
pp=PP+4\{x^{\frac15}+x^{\frac{16}5}+x^{\frac{36}5}+x^{\frac{81}5}+\, .\,.\}
\\[8pt]\quad{}\{x^{\frac45}+x^{\frac95}+x^{\frac{49}5}+x^{\frac{64}5}+\, .\,.\}\\[8pt]
=PP+4xP(x^5,x^4)\, .\,P(x^5,xx)\\[8pt]
=PP+4x\surd\{\tfrac pP\, .\,\tfrac{(Fx^{10})^5}{Fxx}\}\\[8pt]
=PP+4\surd\frac pP\, .\,\frac{\left(\frac{PQR}{2}\right)^{\frac56}}{\left(\frac{pqr}{2}\right)^{\frac16}}
\end{gathered}
\]

\Needspace{30\baselineskip}
\par\medskip\Needspace{5\baselineskip}\pdfbookmark[1]{Art. 10}{transscendentes-v-art-10}\centerline{[10.]}\nobreak\smallskip
Bei der Dreiteilung gilt außerdem, wobei \(p,P\) usw.\ in ihrer bisherigen Bedeutung beibehalten
und \(P(x^{\frac13},1)=P^0\) gesetzt werden
\[
\begin{gathered}
\left(\frac{3PP-P^0P^0}{2}\right)^2=p^4-4\left(\frac{pqr}{2}\right)^{\frac43}\\[8pt]
\left(\frac{3QQ-Q^0Q^0}{2}\right)^2=q^4+4\left(\frac{pqr}{2}\right)^{\frac43}
\end{gathered}
\]
Wenn allgemein \(P(x,1),Q(x,1)\) gegeben sind und \(P(x^n,1),Q(x^n,1)\) gesucht werden,
wobei \(n\) ungerade ist, sind dem \(P(x^n,1)\) die Ausdrücke
\(\pm\surd\frac1n\, .\,P(x^{\frac1n},1)\)
oder \(\pm i\surd\frac1n\, .\,P(x^{\frac1n},1)\) zugeordnet, je nachdem \(n\) von der Form
\(4k+1\) oder \(4k-1\) ist.
Die Gleichung findet sich leicht aus Entwicklung der Summe der geraden Potenzen
der \(n+1\) Wurzeln.

\Needspace{34\baselineskip}
\par\medskip\Needspace{5\baselineskip}\pdfbookmark[1]{Art. 11}{transscendentes-v-art-11}\centerline{[11.]}\nobreak\smallskip
Setzt man
\[
\begin{gathered}P(x^5,y^5)=\frac{R}{r^5}\, .\,P(x,y)^5
\\[8pt]\quad{}+A\, .\,P(x,y)^3Q(x,y)^2
+\frac{5PQR}{pqr^5}\, .\,P(x,y)Q(x,y)^4\end{gathered}
\]

\SourceLeaf{486}
so ist auch
\[
\begin{gathered}
Q(x^5,y^5)
=\frac{5PQR}{pqr^5}\, .\,P(x,y)^4\, .\,Q(x,y)\\[8pt]
\quad+A\, .\,P(x,y)^2\, .\,Q(x,y)^3
+\frac{R}{r^5}Q(x,y)^5
\end{gathered}
\]

\Needspace{25\baselineskip}
Setzt man hier \(y=1\), so erhält man
\[
\begin{gathered}
Ap^3q^3=Pq-\frac{R}{r^5}\, .\,p^5q
-\frac{5PQR}{r^5}\, .\,q^4\\[8pt]
=Qp-\frac{R}{r^5}\, .\,q^5p
-\frac{5PQR}{r^5}\, .\,p^4
\end{gathered}
\]
woraus die frühere Gleichung [Art.\ 9] folgt
\[
0=pQr-Pqr+pqR-5PQR
\]
daraus
\[
A=\frac{r(p^3P-q^3Q)-(p^4+q^4)R}{ppqqr^5}
\]
\clearpage
\endgroup
